\subsection{两个上同伦子的比较}

考虑一个图表

\[
\begin{array}{c} \mathrm{H} \xrightarrow {\varphi} \mathrm{G} \\ u \Big \downarrow \qquad \qquad \qquad \qquad \qquad \qquad \qquad \qquad \qquad \Big \downarrow v \\ \mathrm{H} ^ {\prime} \xrightarrow {\varphi^ {\prime}} \mathrm{G} ^ {\prime} \\ \psi^ {\prime} \end{array}\tag{4}
\]

其中 H, H', G, G' 为群胚，u, v, \(\varphi\) 、\(\psi\) 、\(\varphi'\) 、\(\psi'\) 为群胚态射，满足 \(v \circ \varphi = \varphi' \circ u\) 且 \(v \circ \psi = \psi' \circ u\) 。

以 \(\alpha\) 表示从 G 到上同伦子 \(\mathrm{Coh}(\varphi, \psi)\) 的典范态射，以 \(h\) 表示连接 \(\alpha \circ \varphi\) 与 \(\alpha \circ \psi\) 的典范同伦；类似地定义 \(\alpha'\) 与 \(h'\) 。于是，\(\alpha' \circ v\) 是从 G 到 \(\mathrm{Coh}(\varphi', \psi')\) 的群胚态射，而 \(h' \circ \mathrm{Som}(u)\) 是连接 \(\alpha' \circ \varphi' \circ u\) 与 \(\alpha' \circ \psi' \circ u\) 的同伦，亦即连接 \(\alpha' \circ v \circ \varphi\) 与 \(\alpha' \circ v \circ \psi\) 的同伦。由上同伦子的泛性质（II, 第 185 页, 命题 3），存在唯一的群胚态射 \(w\) ，它从 \(\mathrm{Coh}(\varphi, \psi)\) 到 \(\mathrm{Coh}(\varphi', \psi')\) ，使得 \(w \circ \alpha = \alpha' \circ v\) 且 \(\mathrm{Fl}(w) \circ h = h' \circ \mathrm{Som}(u)\) 。我们特别地把图表 (4) 扩充成了图表

\[
\begin{array}{c} \mathrm{H} \xrightarrow [ \psi ]{\varphi} \mathrm{G} \xrightarrow {\alpha} \mathrm{Coh} (\varphi , \psi) \\ u \Biggl \downarrow \qquad \qquad \qquad \Biggl \downarrow v \qquad \qquad \Biggl \downarrow w \\ \mathrm{H} ^ {\prime} \xrightarrow [ \psi^ {\prime} ]{\varphi^ {\prime}} \mathrm{G} ^ {\prime} \xrightarrow {\alpha^ {\prime}} \mathrm{Coh} (\varphi^ {\prime}, \psi^ {\prime}) \end{array}\tag{5}
\]

其中第二个方块是交换的。

定理 1. —— 假设下列假设成立：

\begin{enumerate}
\def\labelenumi{(\roman{enumi})}
\item
  群胚态射 v 是同伦等价态射；
\item
  映射 \(\operatorname{Orb}(u)\colon\operatorname{Orb}(\mathrm{H})\to\operatorname{Orb}(\mathrm{H}')\) （由 u 通过取轨道而得）是双射；
\item
  在 H 的每个轨道中存在一点 a，使得同态 \(u_a \colon \mathrm{H}_a \to \mathrm{H}_{u(a)}'\) 是满射的。
\end{enumerate}

那么，群胚态射 \(w\colon \operatorname{Coh}(\varphi,\psi)\to\operatorname{Coh}(\varphi',\psi')\) 是同伦等价态射。

设 G'' 为由 G' 经缩并某个极大定向森林的箭而得的群胚。典范态射 \(v'\colon G' \to G''\) 是同伦等价态射（II, 第 184 页, 命题 1 的推论 2）。两个图表

\[
\begin{array}{c c c} \mathrm{H} \xrightarrow [ \psi ]{\varphi} \mathrm{G} & & \mathrm{H} ^ {\prime} \xrightarrow [ \psi^ {\prime} ]{\varphi^ {\prime}} \mathrm{G} ^ {\prime} \\ u \Big \downarrow & \Big \downarrow_ {v ^ {\prime} \circ v} & \text {and} \quad \mathrm{Id} _ {\mathrm{H} ^ {\prime}} \Big \downarrow \\ \mathrm{H} ^ {\prime} \xrightarrow [ v ^ {\prime} \circ \psi^ {\prime} ]{v ^ {\prime} \circ \varphi^ {\prime}} \mathrm{G} ^ {\prime} & & \mathrm{H} ^ {\prime} \xrightarrow [ v ^ {\prime} \circ \psi^ {\prime} ]{v ^ {\prime} \circ \varphi^ {\prime}} \mathrm{G} ^ {\prime \prime} \end{array}
\]

给出群胚态射 \(w_1' \colon \mathrm{Coh}(\varphi, \psi) \to \mathrm{Coh}(v' \circ \varphi', v' \circ \psi')\) 与 \(w_2' \colon \mathrm{Coh}(\varphi', \psi') \to \mathrm{Coh}(v' \circ \varphi', v' \circ \psi')\) ；我们有 \(w_1' = w_2' \circ w\) 。因此只需证明 \(w_1'\) 与 \(w_2'\) 都是同伦等价态射。由于映射 \(\mathrm{Som}(v' \circ v) \colon \mathrm{Som}(\mathrm{G}) \to \mathrm{Som}(\mathrm{G}'' )\) 与 \(\mathrm{Som}(v') \colon \mathrm{Som}(\mathrm{G}') \to \mathrm{Som}(\mathrm{G}'' )\) 都是满射的，从而只需在附加假设——映射 \(\mathrm{Som}(v)\) 是满射的——之下证明定理，我们在余下的证明中始终作此假设。

我们相继证明下列断言：

- 映射 Orb(w) 是双射；

- 对 G 的每个顶点 \(a\) ，同态 \(w_{a}\) 是满射的；

- 对 G 的每个顶点 \(a\) ，同态 \(w_{a}\) 是单射的。

\begin{enumerate}
\def\labelenumi{\alph{enumi})}
\item
  按假设，映射 \(\operatorname{Orb}(u)\) 是双射；由于 v 是同伦等价态射，据 II, 第 182 页, 命题 1，映射 \(\operatorname{Orb}(v)\) 亦然。于是，从偶 \((\varphi,\psi)\) 的框架到偶 \((\varphi',\psi')\) 的框架的、由映射 \(\operatorname{Orb}(u)\) 与 \(\operatorname{Orb}(v)\) 定义的箭图态射是一个同构。特别地，由它通过取连通分支而得的映射是双射。于是 II, 第 185 页, 命题 4 蕴含映射 \(\operatorname{Orb}(w)\) 是双射。
\item
  设 \(f'\) 为 \(G'\) 的箭，以 \(a'\) 表示其起点，\(b'\) 表示其靶。由于映射 \(\operatorname{Som}(v)\) 是满射的，存在 G 中的顶点 \(a\) 与 \(b\) ，使得 \(a' = v(a)\) 且 \(b' = v(b)\) 。由于态射 v 是同伦等价态射，存在 G 的一条箭 \(f\) （它连接 \(a\) 与 \(b\) ），以及元素 \(g \in G_{a}\) ，使得 \(v(g) = f'v(f)^{-1}\) （II, 第 182 页, 命题 1），由此 \(f' = v(gf)\) 。这就证明映射 \(\mathrm{Fl}(v)\) 是满射的。
\end{enumerate}

接着证明映射 \(\mathrm{Fl}(w)\) 是满射的。它的像包含 \(\mathrm{Fl}(\alpha')\) 的像，因为有 \(\alpha'\circ v=w\circ\alpha\) 且映射 \(\mathrm{Fl}(v)\) 是满射的。设 b 为 \(H'\) 的顶点；设 a 为 H 的顶点，使得 b 与 \(u(a)\) 位于 \(H'\) 的同一轨道，且设 f 为 \(H'\) 中连接 \(u(a)\) 与 b 的箭。则，

\[
h ^ {\prime} (u (a)) \cdot (\alpha^ {\prime} \circ \psi^ {\prime}) (f) = (\alpha^ {\prime} \circ \varphi^ {\prime}) (f) \cdot h ^ {\prime} (b),
\]

因为 \(h'\) 是连接 \(\alpha' \circ \varphi'\) 与 \(\alpha' \circ \psi'\) 的同伦。箭 \(h'(u(a)) = w(h(a))\) 属于 \(\mathrm{Fl}(w)\) 的像，由前所述，两条箭 \(\alpha'(\psi'(f))\) 与 \(\alpha'(\varphi'(f))\) 亦然。由此推出箭 \(h'(b)\) 属于 \(\mathrm{Fl}(w)\) 的像，这就证明 \(\mathrm{Fl}(w)\) 的像包含 \(h'\) 的像。由 II, 第 185 页, 命题 2，映射 \(\mathrm{Fl}(w)\) 是满射的。

设 \(g' \in \mathrm{Coh}(\varphi', \psi')_{u(a)}\) 。设 g 为 \(\mathrm{Coh}(\varphi, \psi)\) 的箭，使得 \(w(g) = g'\) 。以 x 与 y 表示 g 的源与靶；我们有 \(u(x) = u(y) = u(a)\) 。设 \(g_1\) （相应地，\(g_2\) ）为 G 中连接 a 与 x（相应地，a 与 y）的箭，其在 v 下的像为 \(e_{u(a)}\) 。则 \(\alpha(g_1)g\alpha(g_2)^{-1}\) 是 \(\mathrm{Coh}(\varphi, \psi)_a\) 的元素，其在 \(w_a\) 下的像为 \(g'\) 。因而，对 G 的每个顶点 a，同态 \(w_a\) 是满射的。

\begin{enumerate}
\def\labelenumi{\alph{enumi})}
\setcounter{enumi}{2}
\tightlist
\item
  我们来证明，对 G 的每个顶点 \(a\) ，同态 \(w_{a}\) 是单射的。相继考虑图表
\end{enumerate}

\[
\begin{array}{c c c} \mathrm{H} \xrightarrow {\varphi} \mathrm{G} & & \mathrm{H} \xrightarrow {v \circ \varphi} \mathrm{G} ^ {\prime} \\ \mathrm{Id} _ {\mathrm{H}} \Biggl \downarrow \quad \psi & \Biggl \downarrow v & u \Biggl \downarrow \quad \Biggl \downarrow \mathrm{Id} _ {\mathrm{G} ^ {\prime}} \\ \mathrm{H} \xrightarrow {v \circ \varphi} \mathrm{G} ^ {\prime} & & \mathrm{H} ^ {\prime} \xrightarrow {\varphi^ {\prime}} \mathrm{G} ^ {\prime} \\ & v \circ \psi & \psi^ {\prime} \end{array}
\]

于是我们归结为处理下列两种情形：1) H' = H 且 \(u = \mathrm{Id}_{\mathrm{H}}\) ；2) G' = G 且 \(v = \mathrm{Id}_{\mathrm{G}}\) 。

\begin{enumerate}
\def\labelenumi{\arabic{enumi})}
\tightlist
\item
  设 H' = H 且 \(u = \mathrm{Id}_{\mathrm{H}}\) 。
\end{enumerate}

考虑一个在同伦意义下为 v 之逆的群胚态射 \(v^{\prime}: G^{\prime} \to G\) ，以及一个同伦 \(k: \operatorname{Som}(G) \to \operatorname{Fl}(G)\) （它连接 \(v^{\prime} \circ v\) 与 \(\mathrm{Id}_{\mathrm{G}}\) ）。由 II, 第 181 页, 注 1，映射 \(\alpha \circ k \circ \varphi\) 与 \(\alpha \circ k \circ \psi\) 分别是连接 \(\alpha \circ v^{\prime} \circ \varphi^{\prime}\) 与 \(\alpha \circ \varphi\) 以及连接 \(\alpha\circ v^{\prime}\circ\psi^{\prime}\) 与 \(\alpha\circ\psi\) 的同伦。因而（参看 II, 第 180 页），映射

\[
\begin{array}{c} h _ {1} \colon \operatorname{Som} (\mathrm{H}) \to \operatorname{Fl} (\operatorname{Coh} (\varphi , \psi)), \\ x \mapsto (\alpha \circ k \circ \varphi) (x) \cdot h (x) \cdot ((\alpha \circ k \circ \psi) (x)) ^ {- 1} \end{array}
\]

是连接 \(\alpha \circ v' \circ \varphi'\) 与 \(\alpha \circ v' \circ \psi'\) 的同伦。由上同伦子的泛性质（II, 第 185 页, 命题 3），存在唯一的群胚态射 \(w'\colon \mathrm{Coh}(\varphi', \psi') \to \mathrm{Coh}(\varphi, \psi)\) ，使得 \(\alpha \circ v' = w' \circ \alpha'\) 且 \(h_1 = \mathrm{Fl}(w') \circ h'\) 。

\[
\begin{array}{c} \mathrm{H} \xrightarrow [ \psi ]{\varphi} \mathrm{G} \xrightarrow {\alpha} \mathrm{Coh} (\varphi , \psi) \\ \mathrm{Id} _ {\mathrm{H}} \Biggl \downarrow \qquad \qquad \qquad \Biggl \downarrow v \qquad \qquad \Biggl \downarrow w \\ \mathrm{H} \xrightarrow [ \psi^ {\prime} ]{\varphi^ {\prime}} \mathrm{G} ^ {\prime} \xrightarrow {\alpha^ {\prime}} \mathrm{Coh} (\varphi^ {\prime}, \psi^ {\prime}) \\ \mathrm{Id} _ {\mathrm{H}} \Biggl \downarrow \qquad \qquad \qquad \Biggl \downarrow v ^ {\prime} \qquad \qquad \Biggl \downarrow w ^ {\prime} \\ \mathrm{H} \xrightarrow [ \psi ]{\varphi} \mathrm{G} \xrightarrow {\alpha} \mathrm{Coh} (\varphi , \psi)  . \end{array}
\]

特别地，我们有

\[
\alpha \circ v ^ {\prime} \circ v = w ^ {\prime} \circ \alpha^ {\prime} \circ v = w ^ {\prime} \circ w \circ \alpha .
\]

由于 k 是连接 \(v'\circ v\) 与 \(\mathrm{Id}_{\mathrm{G}}\) 的同伦，\(\alpha\circ k\) 是连接 \(w'\circ w\circ\alpha\) 与 \(\alpha\) 的同伦。由于 \(\mathrm{Fl}(w'\circ w)\circ h=\mathrm{Fl}(w')\circ h'=h_{1}\) ，故对 H 的每个顶点 x，有

\[
\operatorname{Fl} \left(w ^ {\prime} \circ w\right) \circ h (x) \cdot (\alpha \circ k \circ \psi) (x) = h _ {1} (x) \cdot (\alpha \circ k \circ \psi) (x) = (\alpha \circ k \circ \varphi) (x) \cdot h (x),
\]

此按 \(h_1\) 的定义。把 II, 第 186 页, 命题 5 应用于群胚态射 \(w' \circ w\) 与 \(\mathrm{Id}_{\mathrm{Coh}(\varphi, \psi)}\) ，即知映射 \(\alpha \circ k\) 是连接 \(w' \circ w\) 与 \(\mathrm{Id}_{\mathrm{G}}\) 的同伦。特别地，\(w' \circ w\) 是同伦等价态射。

因而对 G 的每个顶点 \(a\) ，群同态 \((w' \circ w)_a\) 是双射的（II, 第 182 页, 命题 1）。由此推出同态 \(w_a\) 是单射的，从而得到情形 A) 的结论。

\begin{enumerate}
\def\labelenumi{\arabic{enumi})}
\setcounter{enumi}{1}
\tightlist
\item
  设 G' = G 且 \(v = \mathrm{Id}_{\mathrm{G}}\) 。
\end{enumerate}

设 \(x\) 为 H' 的顶点。由于映射 \(\operatorname{Orb}(u)\) 是满射的，存在 H 的顶点 \(a\) 以及 H' 的箭 \(f\) ，它连接 \(u(a)\) 与 \(x\) 。箭 \(\alpha(\varphi'(f))^{-1}\) 、\(h(a)\) 与 \(\alpha(\psi'(f))\) 分别连接 \(\varphi'(x)\) 与 \(\varphi'(u(a)) = \varphi(a)\) 、\(\varphi(a)\) 与 \(\psi(a)\) 以及 \(\psi'(u(a)) = \psi(a)\) 与 \(\psi'(x)\) ；因此它们在 \(\operatorname{Coh}(\varphi, \psi)\) 中可合成。令

\[
h _ {2} (x) = \alpha (\varphi^ {\prime} (f)) ^ {- 1} \cdot h (a) \cdot \alpha (\psi^ {\prime} (f)).
\]

我们来验证这样定义的箭 \(h_{2}(x)\) 不依赖于所选的元素 a 与 f。设 \(a'\) 为 H 的顶点，\(f'\) 为 \(H'\) 中连接 \(u(a')\) 与 x 的箭。由于映射 \(\operatorname{Orb}(u)\) 是单射的，H 的顶点 a 与 \(a'\) 属于同一轨道，且存在箭 \(c \in \operatorname{Fl}(\operatorname{H})\) 连接 a 与 \(a'\) 。于是 \(u(c)f'f^{-1}\) 是 \(u(a)\) 处的圈（在 \(H'\) 中）。由假设 (iii)，存在 H 的顶点 b 使得同态 \(u_{b}\) 是满射的，且存在 H 的从 b 到 a 的箭 \(c'\) ；于是 \(\operatorname{Int}(u(c'))(u(c)f'f^{-1})\) 是 H 中 \(u(b)\) 处的圈，因而它是 \(u_{b}\) 下、b 处的一个圈 \(c''\) 的像。因此，H 的箭 \(g = (c'c)^{-1}c''c'\) 连接顶点 \(a'\) 与顶点 a，且满足 \(f' = u(g)f\) 。于是有

\[
\begin{array}{r l} & {\alpha (\varphi^ {\prime} (f ^ {\prime})) ^ {- 1} h (a ^ {\prime}) \alpha (\psi^ {\prime} (f ^ {\prime}))} \\ & {\qquad = \alpha (\varphi^ {\prime} (f)) ^ {- 1} \alpha (\varphi^ {\prime} (u (g))) ^ {- 1} h (a ^ {\prime}) \alpha (\psi^ {\prime} (u (g))) \alpha (\psi^ {\prime} (f))} \\ & {\qquad = \alpha (\varphi^ {\prime} (f)) ^ {- 1} \alpha (\varphi (g)) ^ {- 1} h (a ^ {\prime}) \alpha (\psi (g)) \alpha (\psi^ {\prime} (f))} \\ & {\qquad = \alpha (\varphi^ {\prime} (f)) ^ {- 1} \cdot h (a) \cdot \alpha (\psi^ {\prime} (f))} \end{array}
\]

因为 h 是连接 \(\alpha \circ \varphi\) 与 \(\alpha \circ \psi\) 的同伦。这就证明了所述的不依赖性。

按构造，我们有 \(h_{2}(u(x)) = h(x)\) （对一切 \(x \in \operatorname{Som}(\mathrm{H})\) ）。我们这样就定义了一个映射 \(h_{2}\) （从 \(\operatorname{Som}(\mathrm{H}')\) 到 \(\operatorname{Fl}(\operatorname{Coh}(\varphi, \psi))\) ）。设 c 为 \(H'\) 的箭，以 x 表示其源，y 表示其靶。设 a 为 H 的顶点，f 为 \(H'\) 中连接 \(u(a)\) 与 x 的箭。则 fc 是 \(H'\) 中连接 \(u(a)\) 与 y 的箭。按 \(h_{2}\) 的定义，我们于是有

\[
\begin{array}{r l} & h _ {2} (x) \alpha (\psi^ {\prime} (c)) = \alpha (\varphi^ {\prime} (f)) ^ {- 1} h (a) \alpha (\psi^ {\prime} (f)) \alpha (\psi^ {\prime} (c)) \\ & \qquad = \alpha (\varphi^ {\prime} (c)) \alpha (\varphi^ {\prime} (f c)) ^ {- 1} h (a) \alpha (\psi^ {\prime} (f c)) \\ & \qquad = \alpha (\varphi^ {\prime} (c)) h _ {2} (y). \end{array}
\]

这就证明 \(h_{2}\) 是连接 \(\alpha \circ \varphi'\) 与 \(\alpha \circ \psi'\) 的同伦。

由上同伦子的泛性质，存在唯一的群胚态射 \(w'\colon \mathrm{Coh}(\varphi', \psi') \to \mathrm{Coh}(\varphi, \psi)\) ，使得 \(w' \circ \alpha' = \alpha\) 且 \(h_2 = \mathrm{Fl}(w') \circ h'\) 。我们有 \(w' \circ w \circ \alpha = w' \circ \alpha' = \alpha\) 与 \(\mathrm{Fl}(w' \circ w) \circ h = \mathrm{Fl}(w') \circ h' \circ \mathrm{Som}(u) = h_2 \circ u = h\) （按 \(h_2\) 的定义）。由上同伦子的泛性质，这蕴含 \(w' \circ w = \mathrm{Id}_{\mathrm{Coh}(\varphi, \psi)}\) 。特别地，对每个 \(a \in \mathrm{Som}(G)\) ，同态 \(w_a\) 是单射的。

于是由 II, 第 182 页, 命题 1 推出，态射 w 是同伦等价态射，定理得证。
% semantic-fragments: legacy-input-seam
\relax{}
